\chapter{Входи машини: як до неї під'єднано очі, вуха та інші датчики}
\chaptermark{Входи машини}
\label{sec:sensors}
\begin{chaptergoals}
\item чому кожне чуття --- перетворювач: рецептор-затвор, регулювання підсилення, кодер імпульсів;
\item чому дробовий шум фотонів обмежує зір у темряві і чому в сутінках ми бачимо повільніше;
\item як регулювання підсилення вміщує величезний діапазон в імпульси, тож датчики передають зміни;
\item як око стискає зображення в сто разів, а вухо працює як банк фільтрів.
\end{chaptergoals}

\section{Датчик як перетворювач}

На рис.~\ref{fig:machine} дані входили в машину зліва, з блока «датчики». Тепер відкриймо цей блок. Для інженера кожен орган чуття --- це \emph{перетворювач} з аналоговим вхідним каскадом і аналого-цифровим перетворювачем наприкінці, тільки цифри на виході --- не числа, а імпульси.

Схема в усіх органів чуття та сама (рис.~\ref{fig:transducer}). Фізична величина --- світло, тиск, коливання, молекула --- діє на білок-рецептор у мембрані чутливої клітини. Цей білок працює як затвор: він сам або через короткий ланцюжок реакцій відкриває іонний канал. Виникає \emph{рецепторний струм}, а з ним --- аналогова напруга на мембрані. Далі напруга проходить регулювання підсилення й потрапляє в кодер імпульсів --- уже знайомий інтегрувальний нейрон~\eqref{eq:glutneuron}, який перетворює рівень на частоту й моменти імпульсів. Вихід майже завжди той самий: чутливі нейрони ока, вуха, нюху й шкіри виділяють глутамат, тож входи під'єднуються просто до шини \nt{GLUT}.

\begin{figure}[t]
\centering
\begin{tikzpicture}[>=Stealth,
  blk/.style={draw,very thick,rounded corners=2pt,align=center,font=\footnotesize,minimum height=1.0cm,text width=1.9cm,fill=blue!5}]
\node[align=center,font=\small] (P) at (0.1,0) {світло, звук,\\тиск,\\молекула};
\node[blk] (R) at (3.0,0) {рецептор-\\затвор};
\node[blk] (G) at (6.5,0) {регулювання\\підсилення};
\node[blk] (E) at (10.0,0) {кодер\\імпульсів};
\node[align=center,font=\small] (B) at (13.4,0) {шина\\\nt{GLUT}};
\draw[->,very thick] (P) -- (R);
\foreach \ic/\xx in {sun/-1.1,speaker/-0.3,press/0.5,molecule/1.3}{\pic[scale=0.85] at (\xx,1.05) {\ic};}
\draw[->,very thick] (R) -- node[above,font=\scriptsize] {струм} (G);
\draw[->,very thick] (G) -- node[above,font=\scriptsize] {рівень} (E);
\draw[->,very thick,red!70!black] (E) -- node[above,font=\scriptsize,black] {імпульси} (B);
\draw[decorate,decoration={brace,amplitude=5pt,mirror},gray!60!black] (1.8,-0.8) -- (7.7,-0.8) node[midway,below=5pt,font=\scriptsize,black] {аналогова частина};
\draw[decorate,decoration={brace,amplitude=5pt,mirror},gray!60!black] (8.8,-0.8) -- (11.2,-0.8) node[midway,below=5pt,font=\scriptsize,black] {імпульси};
\end{tikzpicture}
\caption{Будь-який орган чуття як ланцюг перетворень. Білок-рецептор відкриває канал і дає струм; регулювання підсилення пристосовує його до обставин; кодер~\eqref{eq:glutneuron} перетворює рівень на імпульси, які йдуть на шину \nt{GLUT}. В оці й вусі перші каскади взагалі не видають імпульсів: сигнал лишається аналоговим, доки не доводиться передавати його далеко.}
\label{fig:transducer}
\end{figure}

Одна деталь добре знайома електронникові. В оці й у вусі найперші клітини імпульсів не видають зовсім: вони передають сусідам плавну напругу, як аналоговий каскад на платі. Імпульси з'являються лише там, де сигнал треба везти далеко. Аналоговий сигнал на відстані міліметрів загасає й зашумлюється, а імпульс, як ми бачили в розділі~\ref{sec:neurontypes}, доходить до кінця проводу тієї самої висоти. Оцифровують там, де починається довга лінія.

\section{На межі фізики}

Датчики мозку працюють біля фізичних меж чутливості. Датчик світла в темряві відповідає на \emph{один фотон}: поглинання однієї частинки світла дає струм близько пікоампера, помітний на тлі власного шуму~\cite{Baylor1979}. Датчик звуку помічає зміщення свого чутливого пучка менш ніж на нанометр~\cite{Hudspeth1989}.

Коли світла мало, точність обмежує не датчик, а саме світло: фотони надходять випадково, пуассонівським потоком. Якщо за час накопичення $T$ у датчик потрапляє в середньому $N=\Phi T$ фотонів, їхня кількість тремтить на $\sqrt N$. Щоб приріст $\Delta N$ був помітний, він має в кілька разів перевищувати це тремтіння, і розрізненна частка зміни яскравості дорівнює
\begin{empheq}[box=\keybox]{equation}
\frac{\Delta I}{I} \;\approx\; \frac{k}{\sqrt{\Phi\,T}} .
\label{eq:photonnoise}
\end{empheq}
Це та сама формула, що й для кількості імпульсів~\eqref{eq:rateerror}: дробовий шум фотонів і дробовий шум імпульсів підкоряються одному закону. У темряві око може поліпшити точність лише одним способом --- довше накопичувати фотони, і час накопичення справді зростає до десятих часток секунди. Тому в сутінках ми бачимо повільніше.

\section{Динамічний діапазон: регулювання підсилення}

Найважча інженерна задача для датчиків --- діапазон. Освітленість від зоряної ночі до сонячного снігу змінюється приблизно на десять порядків, інтенсивність звуку від порога чутності до больового --- на дванадцять. А частота імпульсів змінюється від нуля до кількох сотень за секунду, менш ніж на три порядки. Лінійно такий вхід у такий вихід не вкласти.

Рішення те саме, що в будь-якого приймача: \emph{автоматичне регулювання підсилення}. Відповіді чутливих клітин ока добре описує формула
\begin{empheq}[box=\keybox]{equation}
r \;=\; r_{\max}\,\frac{I}{I+\sigma} ,
\label{eq:nakarushton}
\end{empheq}
де $\sigma$ --- яскравість, за якої відповідь сягає половини~\cite{NakaRushton1966}. Сама собою~\eqref{eq:nakarushton} охоплює лише два-три порядки. Але $\sigma$ не стала: за тривалої роботи на яскравому тлі вона зростає слідом за середньою яскравістю. Крива відповіді зсувається вздовж логарифмічної осі яскравості й щоразу ставить свою круту ділянку туди, де зараз перебуває вхід (рис.~\ref{fig:adaptation}). Це те саме ділення на загальну активність, що й шунтувальне гальмування розділу~\ref{sec:gaba}~\cite{Carandini2012}, тільки знаменник тут --- середня яскравість за останні секунди.

\begin{figure}[t]
\centering
\begin{tikzpicture}[>=Stealth,x=1.25cm,y=2.8cm]
\draw[->,gray!70!black] (-2,0) -- (6.4,0) node[below left,font=\small,black] {$\log_{10} I$};
\draw[->,gray!70!black] (-2,0) -- (-2,1.15) node[left,font=\small,black] {$r/r_{\max}$};
\foreach \u in {-2,0,2,4,6}{\draw[gray!60!black] (\u,-0.02) -- (\u,0.02); \node[below,font=\scriptsize] at (\u,0) {$\u$};}
\draw[densely dashed,gray!60!black] (-2,0.5) -- (6.2,0.5);
\foreach \s/\c/\lab in {0/blue!60!black/{темне тло},2/green!45!black/{середнє},4/red!70!black/{яскраве}}{
  \pgfmathsetmacro{\lo}{max(-2,\s-3)}
  \draw[very thick,\c] (-2,0) -- (\lo,0) plot[domain=\lo:6.2,samples=80] (\x,{1/(1+10^(\s-\x))});
  \node[font=\scriptsize,\c,anchor=west] at (\s+0.15,0.42) {\lab};}
\foreach \s/\ic in {0/moon,2/cloud,4/sun}{\pic at (\s+0.6,0.64) {\ic};}
\end{tikzpicture}
\caption{Регулювання підсилення датчика світла за~\eqref{eq:nakarushton}. Кожна крива охоплює два-три порядки яскравості; при переході на яскравіше тло $\sigma$ зростає, і крива зсувається логарифмічною віссю. Разом криві, що зсуваються, покривають увесь діапазон.}
\label{fig:adaptation}
\end{figure}

Регулювання має ціну, знайому з розділу~\ref{sec:code}: датчик повідомляє не яскравість, а яскравість \emph{відносно тла}. Сталий вхід поступово перестає викликати відповідь, а кожна зміна викликає. Входи машини від самого початку говорять про зміни, а не про рівні, і в цьому та сама економія імпульсів, що й у твердженні~\ref{prop:economy}~\cite{Barlow1961}. Звідси й датчики увімкнення та вимкнення розділу~\ref{sec:glutcomputer}~\cite{Schiller1992}: одні повідомляють, що стало яскравіше, інші --- що темніше.

Інженери повторили цей прийом у \emph{подієвих камерах}. Така камера не знімає кадрів за тактовим сигналом: кожен її піксель сам, без спільного такту, видає подію «яскравіше» або «темніше», коли логарифм яскравості змінився на задану частку~\cite{Lichtsteiner2008}. Це точнісінько двопровідний код увімкнення---вимкнення з розділу~\ref{sec:glutcomputer}, втілений у кремнії, і він дає камері діапазон і швидкодію, недосяжні для звичайної матриці~\cite{Gallego2022}.

\section{Око: стиснення до проводу}

В оці близько $10^8$ датчиків світла~\cite{Curcio1990}, а вихідних нейронів, які везуть зображення в мозок, --- близько мільйона. Перш ніж сигнал залишить око, його стискають приблизно в сто разів. Головні прийоми стиснення нам знайомі. Кожен вихідний нейрон відповідає на різницю між яскравістю в маленькій плямі та яскравістю навколо неї --- це віднімання сусідів через гальмування, як у розділі~\ref{sec:gaba}. Рівні ділянки зображення майже нічого не коштують, а краї й зміни передаються. Різні вихідні нейрони спеціалізовані: одні повідомляють про посвітлішання, інші про потемніння, ще інші про рух у певний бік (рис.~\ref{fig:direction}); у миші таких вихідних каналів налічили понад тридцять~\cite{Baden2016}.

Яка пропускна здатність отриманого кабелю? У морської свинки за записами вихідних нейронів виміряли близько $875$ тисяч біт за секунду на $\sim10^5$ таких нейронів; перерахунок на мільйон вихідних нейронів людини дає порядку $10^7$ біт за секунду~\cite{Koch2006} --- стільки ж, скільки старий мережевий кабель на десять мегабіт. За середньої частоти в кілька імпульсів за секунду на нейрон це кілька бітів на імпульс, як і виходило в розділі~\ref{sec:code}.

\section{Вухо: банк фільтрів}

Вухо розв'язує іншу задачу: звук треба розкласти за частотами. Інженер поставив би банк смугових фільтрів, і вухо робить саме це механічно. Звукове коливання біжить пружною перегородкою, властивості якої плавно змінюються вздовж її довжини, і кожна ділянка резонує на свою частоту. Датчики, що сидять на ділянці, відповідають лише на свою смугу (рис.~\ref{fig:filterbank}). Частоти розкладено за місцем приблизно логарифмічно: кожній октаві дістається схожа частка датчиків. Номер датчика, що спрацював, --- це частота, код місцем із розділу~\ref{sec:code}.

\begin{figure}[t]
\centering
\begin{tikzpicture}[>=Stealth,x=1.35cm,y=2.2cm]
\draw[->,gray!70!black] (-0.3,0) -- (7.6,0) node[above left,font=\small,black] {частота, Гц};
\draw[->,gray!70!black] (-0.3,0) -- (-0.3,1.15) node[left,font=\small,black] {відповідь};
\foreach \k/\f in {0/125,1/250,2/500,3/1000,4/2000,5/4000,6/8000}{
  \draw[gray!60!black] (\k,-0.02) -- (\k,0.02); \node[below,font=\scriptsize] at (\k,0) {\f};
  \draw[thick,blue!60!black] plot[domain={max(\k-1.2,-0.3)}:{\k+0.7},samples=60] (\x,{1/(1+((\x-\k)/(\x<\k?0.45:0.2))^4)});}
% a piano keyboard: seven white keys per octave, like the octaves on the axis
\foreach \i in {0,...,48}{\draw[gray!50!black,fill=white,line width=0.3pt] ({\i/7},-0.44) rectangle ({(\i+1)/7},-0.26);}
\foreach \o in {0,...,6}{\foreach \j in {0,1,3,4,5}{\fill[black] ({\o+(\j+1)/7-0.035},-0.35) rectangle ({\o+(\j+1)/7+0.035},-0.26);}}
\end{tikzpicture}
\caption{Вухо як банк смугових фільтрів, схема. Кожна крива --- відповідь датчиків однієї ділянки на звук різної частоти; центральні частоти йдуть через октаву, як клавіші на логарифмічній шкалі. У справжніх фільтрів крутий високочастотний схил і пологий низькочастотний. Унизу --- клавіатура: на ній, як і у вусі, кожній октаві дістається однакове місце.}
\label{fig:filterbank}
\end{figure}

Резонатори вуха не пасивні. Частина датчиків сама розгойдує перегородку в такт звуку й підсилює слабкі коливання в тисячі разів за амплітудою (на $66$--$76$\,дБ), а зі зростанням гучності підсилення падає~\cite{Ruggero1997,Robles2001}. Для інженера це підсилювач із додатним зворотним зв'язком, поставлений біля самої межі самозбудження, --- те саме життя на межі, що й у мережі розділу~\ref{sec:gaba}: біля межі система особливо чутлива до малих сигналів. Стиснення гучності тут робить той самий підсилювач: тихе підсилюється сильно, гучне --- слабко.

У вуха є й другий код. На низьких частотах імпульси нейронів ідуть у фазі зі звуком: нейрон спрацьовує на певній ділянці кожної хвилі, хоч і не на кожній хвилі. У морської свинки прив'язка до фази починає слабшати близько $600$\,Гц і ще помітна приблизно до $3.5$\,кГц~\cite{PalmerRussell1986}. Так у моментах імпульсів зберігається точний час звуку, а з нього --- різниця приходу звуку у два вуха, за якою мережа розділу~\ref{sec:glutcomputer} знаходить напрямок~\cite{Grothe2010}. На високих частотах імпульси за хвилею не встигають, і лишається тільки код місцем. Одне вухо користується обома кодами розділу~\ref{sec:code}: часом там, де нейрони встигають, і місцем там, де ні.

\section{Інші датчики}

\begin{table}[t]
\centering
\footnotesize
\setlength{\tabcolsep}{3pt}
\renewcommand{\arraystretch}{1.15}
\begin{tabular}{>{\raggedright}p{1.9cm}>{\raggedright}p{2.6cm}>{\raggedright}p{4.0cm}>{\raggedright\arraybackslash}p{4.6cm}}
\toprule
Чуття & Що вимірює & Як улаштовано вхід & Прийом із книги \\
\midrule
зір & світло & $\sim10^8$ датчиків, стиснення в $\sim100$ разів, $\sim10^7$ біт/с & регулювання підсилення, віднімання сусідів, два проводи, напрямок руху \\
слух & тиск повітря & механічний банк фільтрів, активний підсилювач & код місцем і код часом, затримки \\
дотик & тиск, вібрація & датчики, що звикають швидко й повільно & пара «фільтр верхніх частот --- фільтр нижніх» \\
температура & тепло & окремі датчики тепла й холоду & двопровідний код \\
нюх & молекули & $\sim400$ типів рецепторів, у кожного датчика один тип & комбінаторний код: запах --- набір типів, що спрацювали \\
смак & речовини в їжі & п'ять якостей: солодке, гірке, солоне, кисле, умамі & мічені лінії; частина клітин виділяє АТФ або серотонін, а не глутамат \\
положення тіла & розтяг м'язів & датчики довжини й швидкості розтягу & зворотний зв'язок для регулятора руху \\
\bottomrule
\end{tabular}
\caption{Як під'єднано датчики. Усі пристрої в третьому стовпці виміряно; правий стовпець пов'язує їх із прийомами з попередніх розділів.}
\label{tab:sensors}
\end{table}

Інші чуття зібрано в таблиці~\ref{tab:sensors}, і майже в кожному рядку впізнається прийом із попередніх розділів. Дотик користується парою фільтрів: одні датчики тиску відповідають на дотик і відразу замовкають, інші тримають відповідь, доки тиск є. Температуру передають два проводи, окремі для тепла й для холоду. Найнезвичайніший вхід --- нюх. Типів нюхових рецепторів у людини близько чотирьохсот, і кожен нюховий нейрон несе лише один тип~\cite{Mombaerts2004}. Запах активує не один тип, а набір, і повідомлення --- це \emph{які} типи спрацювали. Для інженера це двійковий вектор довжиною близько чотирьохсот, кількість можливих значень якого астрономічна.

\begin{keyblock}
\begin{proposition}[Входи машини]
\label{prop:sensors}
Кожне чуття --- перетворювач, який працює біля фізичної межі чутливості, стискає величезний динамічний діапазон автоматичним регулюванням підсилення й передає шиною \nt{GLUT} насамперед \emph{зміни}. Для цього датчики користуються тими самими прийомами, що й сама машина: двопровідним кодом, кодом місцем, кодом часом і діленням на тло. Будову датчиків виміряно; те, що їхні коди підлаштовано під найбільшу інформацію на імпульс, --- добре обґрунтована гіпотеза.
\end{proposition}
\end{keyblock}

Входи, як і все інше, працюють асинхронно. Кожен датчик видає імпульс, коли йому є що повідомити, а різні чуття надходять із різними затримками: звук --- через мілісекунди, світло --- через десятки мілісекунд. Як машина зводить їх в одну картину світу, якщо спільного такту немає, --- одна із задач узгодження в часі з розділу~\ref{sec:synthesis}.

\section{Задачі}

\begin{exercise}[\lvlA]
\label{ex:11:1}
\emph{Скільки накопичувати фотони.} У сутінках у датчик потрапляє $100$ фотонів за секунду; приріст помітний, якщо він у $3$ рази більший за тремтіння~\eqref{eq:photonnoise}. (а)~Скільки часу потрібно одному датчикові, щоб помітити зміну яскравості на $10\,\%$? (б)~Скільки датчиків треба підсумувати, щоб устигнути за $0.2\,\text{с}$, і в скільки разів упаде роздільність за віссю?
\end{exercise}

\begin{exercise}[\lvlA]
\label{ex:11:2}
\emph{Скільки бітів в імпульсі.} (а)~Око передає $\sim10^7$ біт/с через $10^6$ вихідних нейронів із частотою $3$--$5$ імпульсів за секунду. Яку частку пропускної здатності~\eqref{eq:spikeentropy} при $\Delta t=1\ms$ воно використовує? (б)~Запах вмикає $20$ з $\approx400$ типів рецепторів. Скільки бітів несе такий набір і скільки спільних типів у середньому мають два випадкові запахи?
\end{exercise}

\begin{exercise}[\lvlB]
\label{ex:11:3}
\emph{Що вміє одна крива~\eqref{eq:nakarushton}.} (а)~У якому діапазоні яскравостей відповідь зростає від $10\,\%$ до $90\,\%$ від $r_{\max}$? Скільки це порядків? (б)~Скільки положень зсувної кривої потрібно, щоб покрити десять порядків яскравості? дванадцять порядків гучності?
\end{exercise}

\begin{exercise}[\lvlB]
\label{ex:11:4}
\emph{Межа коду часом у вусі.} Різниця приходу звуку у два вуха --- до $0.22\,\text{м}/343\,\text{м/с}$. (а)~Імпульси йдуть у фазі з тоном: до якої частоти напрямок за ними визначається однозначно? (б)~Імпульс тремтить на $0.5\ms$, а розрізняти треба $20$~мкс. Скільки нейронів по $100$ імпульсів за секунду треба усереднити за $50\ms$?
\end{exercise}

\begin{exercise}[\lvlB]
\label{ex:11:5}
\emph{Подієва камера на кабель ока.} Камера з $10^6$ пікселів із виходом $10^7$ біт/с надсилає подію (адресу й знак), коли $\ln I$ у пікселі змінився на $\ln1.2$. З якою найбільшою швидкістю може бігти край довжиною $500$ пікселів із перепадом яскравості $4$? Скільки кадрів за секунду дала б через той самий вихід звичайна камера по $8$ бітів на піксель?
\end{exercise}

\begin{exercise}[\lvlB]
\label{ex:11:6}
\emph{Моделювання: регулювання підсилення.} Датчик~\eqref{eq:nakarushton} підлаштовує $\sigma$ з $\tau=1\,\text{с}$ у логарифмах, $\tau\,\dd\ln\sigma/\dd t=\ln I-\ln\sigma$, або лінійно, $\tau\,\dd\sigma/\dd t=I-\sigma$; тло стрибає $1\to100\to1$. Через скільки секунд відповідь на пробу $+10\,\%$ відновлюється до половини звиклої після стрибка вгору і вниз за кожного закону?
\end{exercise}

\begin{exercise}[\lvlC]
\label{ex:11:7}
\emph{Під який світ підлаштовано криву.} За малого сталого шуму виходу інформація про яскравість найбільша, коли $r(I)=r_{\max}\Phi_I(I)$, де $\Phi_I$ --- функція розподілу яскравостей. (а)~Для якого розподілу оптимальна крива~\eqref{eq:nakarushton} і який його розкид $\log_{10}I$? (б)~Яка крива оптимальна за пуассонівського шуму виходу?
\end{exercise}
